## TEMP holding note — RQ3 category attribution, RERUN on current VIF-screened Set4 (2026-08-11)

**Supersedes** `TEMP_rq3_category_attribution_results_2026-08-10.tex` (built on the pre-VIF
`nested_set4_gated_all`, 25 columns — kept as historical record, not deleted). This rerun uses
the current, VIF-screened `nested_set4_gated_all_vif` (20 columns) — same script
(`run_rq3_category_attribution.py`), same method (`grouped_permutation_importance` +
`category_morans_i_before_after`), single spatial_block split, 4survey only. Ran locally
2026-08-11.

### Results — permutation importance by category

| Category | n columns | mean R2 drop (shuffled) | 2026-08-10 (pre-VIF) equivalent |
|---|--:|---:|---:|
| stand_structure | 4 | 0.2949 | 0.3072 |
| terrain | 3 | 0.0894 | 0.1176 |
| climate | 2 | **0.0640** | 0.0297 |
| spatial_position_edge_effects | 3 | 0.0579 | 0.0360 |
| wind | 3 | 0.0534 | 0.0334 |
| soil_site | 5 | **-0.0007** | 0.0020 |

### Results — Moran's I before/after removing each category (baseline full-model R2 = 0.1501, Moran's I = 0.1525)

| Category removed | R2 | Moran's I |
|---|---:|---:|
| terrain | 0.0968 | 0.0417 |
| wind | 0.1478 | 0.1382 |
| soil_site | 0.1250 | 0.0245 |
| climate | 0.1292 | 0.1236 |
| spatial_position_edge_effects | 0.1485 | 0.0422 |
| stand_structure | 0.0025 | 0.0220 |

### Headline finding — a real update, not just a refresh

**`stand_structure` still dominates by a wide margin** (0.295, ~3.3x the next-largest category) —
this part of the earlier finding holds up under VIF screening.

**But the earlier "every other category removal IMPROVES R2" pattern is GONE.** On the pre-VIF
Set4 (2026-08-10), removing terrain/wind/soil/climate/edge each *increased* R2 above the
full-model baseline — the interpretation then was that most environmental columns were adding
noise, not signal, on that single split. On the VIF-screened Set4, removing terrain/soil/climate
now *decreases* R2 (0.150 → 0.097/0.125/0.129), and removing wind/edge barely changes it (0.150 →
0.148/0.148) — **no category removal improves the model anymore**. This is consistent with VIF
screening doing exactly its intended job: the pre-VIF "noise" was likely collinear/redundant
columns competing with each other and adding variance without adding independent signal; removing
that collinearity lets terrain/climate/soil contribute real, if still secondary, predictive value.

**`soil_site` remains the weakest category, now indistinguishable from pure noise** (-0.0007,
essentially zero, on 5 columns) — consistent across both the pre-VIF and VIF-screened versions,
the one part of the original finding that hasn't changed at all.

**Climate's contribution roughly doubled** (0.030 → 0.064) — worth citing as real, not an
artifact, since VIF screening specifically removed collinear climate-adjacent columns that were
previously diluting its apparent individual contribution.

**Moran's I pattern still doesn't cleanly fit the "removing a real category raises spatial
clustering" framework** the check was designed around — removing terrain, the category with the
largest real R2 contribution besides stand_structure, actually *lowers* residual spatial
clustering (0.153 → 0.042) rather than raising it. This was true in the pre-VIF version too and
remains unexplained — flagged as a genuine open question, not resolved by this rerun, rather than
force-fit an interpretation.

### Still applies from the original note

The confound/circularity check on `CanopyCover` (how it's actually derived) is still not done —
`stand_structure`'s dominance should still be read with that caveat pending.
